% Sucesor del articulo V: incorporacion focal F051.
% Fragmento bilingue emparejado con la homologa EN; no compilar aisladamente.
\subsection{Observador primo de vacancias y resoluci\'on modal}
\label{sec:f051-prime-vacancy-observer}

En la cadena causal del estado HMT ya generado, el producto non\'adico, su
coproducto multiplicativo reducido y la proyecci\'on espectral
\(P_{\rm irr}\) producen las formas irreducibles; la evaluaci\'on
\(\operatorname{ev}_9\) las reconoce despu\'es como primos ordinarios. La
frontera \(729/1000\) produce el calendario de vacancias. El observador
\(\mathcal C_{\varepsilon,X}\) compone ambos objetos internos con el reloj
logar\'itmico. La criba convencional interviene únicamente como lector
externo de la reproducci\'on num\'erica.

\subsubsection{Calendario y equivalencia de las dos escrituras}

Sea
\begin{equation}
 \varepsilon_{\rm vac}
 =1-\log_{10}9
 =\log_{10}\!\left(\frac{10}{9}\right),
 \qquad
 z_n=\left\lceil(n+1)\varepsilon_{\rm vac}\right\rceil
      -\left\lceil n\varepsilon_{\rm vac}\right\rceil ,
 \quad n\geq1.
 \label{eq:f051-epsilon-calendar}
\end{equation}
La cantidad \(\varepsilon_{\rm vac}\) es irracional. En efecto, si
\(\varepsilon_{\rm vac}=a/b\in\mathbb Q\), entonces
\((10/9)^b=10^a\), lo que contradice la unicidad de la factorizaci\'on
prima por la presencia de una potencia no nula de \(3\). En consecuencia,
ni \(n\varepsilon_{\rm vac}\) ni
\((n+1)\varepsilon_{\rm vac}\) son enteros para \(n\geq1\), y
\begin{equation}
 z_n=\left\lfloor(n+1)\varepsilon_{\rm vac}\right\rfloor
      -\left\lfloor n\varepsilon_{\rm vac}\right\rfloor
 \in\{0,1\}.
 \label{eq:f051-ceil-floor}
\end{equation}
La restricci\'on \(n\geq1\) es necesaria para esta equivalencia: en
\(n=0\), la convenci\'on con techos y la convenci\'on con suelos no coinciden.

\subsubsection{Cadena tipada y observador de correlaci\'on}

Para \(X\geq2\), la cadena causal puede escribirse sin p\'erdida de tipos
como
\begin{equation}
 \operatorname{Ran}(P_{\rm irr})
 \xrightarrow{\ \operatorname{ev}_9\ }
 \mathbb P
 \xrightarrow{\ \iota\ }
 \{0,1\}\times\mathbb R_{>0},
 \qquad
 \bigl((z_p,\log p)\bigr)_{p\in\mathbb P_{\leq X}}
 \xrightarrow{\ \mathcal C_{\varepsilon,X}\ }
 \mathbb R,
 \label{eq:f051-typed-chain}
\end{equation}
donde
\[
 \iota(p)=(z_p,\log p),
 \qquad
 \mathcal C_{\varepsilon,X}\colon
 \prod_{p\in\mathbb P_{\leq X}}
       (\{0,1\}\times\mathbb R_{>0})\longrightarrow\mathbb R,
 \qquad
 \mathcal C_{\varepsilon,X}\bigl((a_p,\ell_p)_{p\leq X}\bigr)
 =\sum_{p\leq X}(a_p-\varepsilon_{\rm vac})\ell_p .
\]
Por tanto,
\begin{equation}
 \theta_{\rm vac}(X)=\sum_{p\leq X}z_p\log p,
 \qquad
 \theta(X)=\sum_{p\leq X}\log p,
 \qquad
 R_{\rm vac}(X)=\theta_{\rm vac}(X)
                   -\varepsilon_{\rm vac}\theta(X).
 \label{eq:f051-rvac-chebyshev}
\end{equation}
El centrado exacto del observador es
\(\varepsilon_{\rm vac}\theta(X)\); la expresi\'on
\(\varepsilon_{\rm vac}X\) corresponde a otra normalizaci\'on. Dos prefijos
con la misma cantidad de vacancias pueden producir valores distintos de
\(R_{\rm vac}\), porque el observador conserva la posici\'on prima y su peso
logar\'itmico.

\subsubsection{Lectores de Dirichlet y dominios analíticos}

El observador finito y los dos lectores de Dirichlet constituyen
realizaciones distintas y compatibles:
\begin{align}
 Z^{+}_{\rm vac}(s)
   &=\sum_{n\geq1}\frac{z_n}{n^s}
     =\varepsilon_{\rm vac}\zeta(s)+H_{\rm vac}(s),
 & H_{\rm vac}&\in\mathcal O(\operatorname{Re}s>0),
 \nonumber\\
 Z^{\times}_{\rm vac}(s)
   &=\prod_p(1-p^{-s})^{-z_p}
     =\zeta(s)^{\varepsilon_{\rm vac}}G_{\rm vac}(s),
 & \operatorname{Re}s&>1,
 \label{eq:f051-dirichlet-readers}\\
 \log G_{\rm vac}(s)
   &=\sum_p(z_p-\varepsilon_{\rm vac})
       \log(1-p^{-s})^{-1},
 & \operatorname{Re}s&>1.
 \nonumber
\end{align}
La primera serie y el producto de Euler convergen absolutamente para
\(\operatorname{Re}s>1\). La discrepancia parcial de la palabra mec\'anica
es acotada, y la sumaci\'on de Abel prolonga \(H_{\rm vac}\) de forma
holomorfa a \(\operatorname{Re}s>0\). La potencia
\(\zeta(s)^{\varepsilon_{\rm vac}}\) se define aqu\'i mediante el logaritmo
de Euler en el semiplano de convergencia absoluta.

\subsubsection{Resoluci\'on de Fourier y realidad del observador}

Def\'inanse, para \(k\in\mathbb Z\setminus\{0\}\),
\begin{equation}
 c_k=\frac{e^{2\pi i k\varepsilon_{\rm vac}}-1}{2\pi i k},
 \qquad
 S_k(X)=\sum_{p\leq X}(\log p)
             e^{2\pi i k p\varepsilon_{\rm vac}}.
 \label{eq:f051-fourier-data}
\end{equation}

\begin{proposition}[Descomposici\'on modal de la correlaci\'on prima]
\label{prop:f051-fourier-resolution}
Para todo \(X\geq2\),
\begin{equation}
 R_{\rm vac}(X)
 =\lim_{K\to\infty}
   \sum_{0<|k|\leq K}c_kS_k(X)
 =\lim_{K\to\infty}
   \sum_{0<|k|\leq K}
   \left(1-\frac{|k|}{K+1}\right)c_kS_k(X).
 \label{eq:f051-fourier-pair}
\end{equation}
El primer l\'imite se define mediante sumas parciales sim\'etricas y el
segundo mediante medias de Fej\'er; éstas son exactamente las dos nociones
de convergencia afirmadas. La convergencia absoluta en \(k\) queda fuera
de las hip\'otesis y de la conclusi\'on.
\end{proposition}

\begin{proof}
Sea
\[
 f_{\varepsilon}(x)
 =\mathbf1_{(1-\varepsilon_{\rm vac},1)}(\{x\})
  -\varepsilon_{\rm vac}.
\]
Sus coeficientes de Fourier no nulos son precisamente \(c_k\). La
irracionalidad de \(\varepsilon_{\rm vac}\) impide que
\(p\varepsilon_{\rm vac}\) alcance un extremo del intervalo; por
\eqref{eq:f051-ceil-floor},
\(f_{\varepsilon}(p\varepsilon_{\rm vac})
 =z_p-\varepsilon_{\rm vac}\). El teorema de Dirichlet da la convergencia
de las sumas sim\'etricas en esos puntos de continuidad, y el teorema de
Fej\'er da la convergencia de las medias. La finitud del conjunto
\(p\leq X\) permite intercambiar el l\'imite con la suma prima ponderada en
los sentidos sim\'etrico y de Fej\'er. La estimaci\'on
\(|c_k|=O(|k|^{-1})\) no implica sumabilidad absoluta y, por tanto, no
habilita una reordenaci\'on absoluta.

Finalmente,
\(c_{-k}=\overline{c_k}\) y
\(S_{-k}(X)=\overline{S_k(X)}\). Los modos \(k\) y \(-k\) contribuyen
\(2\operatorname{Re}(c_kS_k(X))\); de aqu\'i se recupera directamente la
realidad de \(R_{\rm vac}(X)\).
\end{proof}

\subsubsection{Testigos numéricos y reproducción independiente}

La Tabla~\ref{tab:f051-rvac-powers} presenta los seis cortes finitos
disponibles. Los valores de \(R_{\rm vac}\) reproducen todos los d\'igitos
almacenados y la columna normalizada conserva el redondeo de la fuente
num\'erica.

\begin{table}[htbp]
\centering
\small
\begin{tabular}{rrrrr}
\toprule
\(X\) & \(\pi(X)\) & \(\pi_{\rm vac}(X)\) & \(R_{\rm vac}(X)\)
 & \(R_{\rm vac}(X)/\sqrt X\)\\
\midrule
\(10^3\) & 168 & 7 & \(-4.553590011317908\) & \(-0.143997159663565\)\\
\(10^4\) & 1\,229 & 61 & \(41.334631595982586\) & \(0.413346315959826\)\\
\(10^5\) & 9\,592 & 427 & \(-141.1579619980301\) & \(-0.446380669781268\)\\
\(10^6\) & 78\,498 & 3\,595 & \(32.21713602110386\) & \(0.032217136021104\)\\
\(10^7\) & 664\,579 & 30\,384 & \(-418.47581825715594\) & \(-0.132333673139529\)\\
\(10^8\) & 5\,761\,455 & 263\,841 & \(4021.0180144347314\) & \(0.402101801443473\)\\
\bottomrule
\end{tabular}
\caption{Observador de vacancias sobre primos en seis cortes finitos.}
\label{tab:f051-rvac-powers}
\end{table}

La tabla completa de treinta modos reside en
\texttt{S\_k\_1\_30\_X1e8.csv}. Sus cinco mayores valores de
\(|S_k(10^8)|\), ordenados por m\'odulo, son:

\begin{table}[htbp]
\centering
\scriptsize
\setlength{\tabcolsep}{3.5pt}
\begin{tabular}{rrrrrr}
\toprule
\(k\) & \(\operatorname{Re}S_k\) & \(\operatorname{Im}S_k\)
 & \(|S_k|\) & \(|S_k|/\sqrt X\) & \(|S_k|/(\sqrt X\log X)\)\\
\midrule
19 & 30646.583055876723 & -39786.84285255695 & 50221.56824686792 & 5.0221568246867925 & 0.2726368745267788\\
24 & -1325.0060816425494 & 47631.087438932576 & 47649.51344695589 & 4.764951344695589 & 0.25867400944234675\\
21 & 1214.7582244307505 & -44824.22356680968 & 44840.68081453668 & 4.484068081453668 & 0.24342575303172864\\
26 & -26879.47488364538 & 33364.6769815865 & 42845.161221614406 & 4.284516122161441 & 0.23259271368502904\\
1 & 40391.94841735065 & 13384.21089257329 & 42551.693246765084 & 4.255169324676508 & 0.23099956965887428\\
\bottomrule
\end{tabular}
\caption{Cinco mayores modos en la muestra \(1\leq k\leq30\) para
\(X=10^8\). En este dominio muestral,
\(\lvert S_{19}(10^8)\rvert\) es el máximo.}
\label{tab:f051-top-five-modes}
\end{table}

El cálculo independiente de las cuatro primeras filas, obtenido desde
\eqref{eq:f051-epsilon-calendar} con una criba usada como lector de
verificaci\'on, reproduce los valores hasta \(10^6\) con errores absolutos
inferiores a \(2.1\times10^{-11}\). Las filas de \(10^7\) y \(10^8\) y los
treinta modos son testigos num\'ericos finitos del c\'alculo hist\'orico. La
proyecci\'on \(P_{\rm irr}\) de \eqref{eq:f051-typed-chain} conserva el papel
de generador interno.

\subsubsection{Condición analítica para la cota crítica}

La estimaci\'on
\begin{equation}
 R_{\rm vac}(X)=O\!\left(X^{1/2}\log^B X\right)
 \label{eq:f051-conditional-bound}
\end{equation}
controlar\'ia el t\'ermino primo centrado del factor \(G_{\rm vac}\). Su
demostraci\'on exige cotas para \(S_k(X)\) uniformes en \(k\) y compatibles
con la suma sim\'etrica o de Fej\'er. Las tablas caracterizan una ventana
finita y no aportan esa uniformidad. El resultado demostrado es la
identidad modal exacta para cada \(X\) finito; la cota asint\'otica cr\'itica
permanece condicionada a las estimaciones modales uniformes.
